\section{Conclusions}
\label{sec:conclusions}

This systematic literature review has thoroughly examined and synthesized the existing body of research regarding AI techniques applied to customer churn prediction within the retail sector, with particular attention to non-contractual environments characterized by gradual customer disengagement. The increasing digitalization of retail and higher competition make it easier for customers to defect and switch to competitors. Hence, understanding customer churn has become critical for businesses seeking strategic customer retention.
%Our investigation addressed both publication space and research space questions. 
Our investigation identified a notable increase in research interest from 2021 onwards, highlighting a timely expansion of academic attention likely driven by the growing demand for such research by the industry.

The wide range of research areas covered by the 41 analyzed studies reflects a clear interdisciplinary approach, driven by the complexity and multifaceted nature of churn prediction problems.
%Geographically, China, India, Belgium, and Iran emerge as key contributors to this field, reflecting the global significance and applicability of churn prediction research. However, cross-national collaborations remain infrequent, suggesting potential opportunities for increased international research synergy and knowledge exchange.
In terms of churn definition, considerable variability was observed. The predominant approach adopted was based on a fixed time window of inactivity, followed closely by dynamic definitions. Each approach offers distinct advantages, from simplicity and ease of application to greater sensitivity in capturing subtle shifts in customer behavior. Nonetheless, the diversity of definitions highlights an ongoing challenge in the field: establishing standardized, universally applicable criteria for churn identification, which is fundamental for consistent comparison and evaluation of predictive models.

Regarding features for churn prediction, transactional features were found to be central, as they directly reflect customer interactions with retailers. Behavioral and demographic data were also widely employed, reiterating the importance of such data in capturing complex customer behaviors. Feature engineering techniques, particularly RFM analyses, were among the most effective and widely used methods, highlighting the importance of intelligent feature transformation to improve prediction performance.

The analysis of AI techniques revealed a clear preference for and effectiveness of tree-based algorithms, particularly Random Forest and XGBoost, as well as neural network approaches, including Long Short-Term Memory networks. More recent and advanced techniques, such as time-series analysis and Transformer-based models, have yet to be widely explored, despite their potential to capture complex patterns and temporal dependencies in customer behavior.

Model performance evaluation was predominantly conducted using Accuracy, Precision, Recall, F1-score, and AUC-ROC metrics, highlighting the importance of accurately identifying churners in scenarios where false positives and negatives bear significant operational costs. 

However, several challenges persist. Notably, the limited availability of sufficiently large and diverse real-world datasets poses substantial constraints on model robustness and generalizability. Moreover, the computational cost and scalability of sophisticated AI models remain significant barriers, necessitating optimized computational strategies and the adoption of scalable cloud-based infrastructure. When studies compare different AI techniques, the observed performance may not reflect each method's full potential, as resource constraints can limit the ability to optimize hyperparameters effectively for each method.

Emerging trends from this review clearly indicate opportunities for future research to expand the depth and breadth of churn prediction methodologies. Particularly promising areas include the integration of richer and more diverse data sources, such as social media %sentiment analysis
data, detailed transaction histories, and socio-economic indicators, which can significantly enhance predictive accuracy and insights. Furthermore, advancing explainable AI models that leverage intrinsic attention mechanisms is strongly recommended to bridge the gap between model complexity and interpretability, enabling clearer strategic guidance for practitioners.

Additionally, the studies highlight key strategies for enhancing the predictive performance of churn models. These include the use of automated hyperparameter tuning to optimize model performance, the application of advanced ensemble methods, and the adoption of modern AI architectures such as transformer-based models and recurrent neural networks. Exploring spatial dynamics in churn behavior---especially the transition from partial to complete churn---also represents a compelling direction for future investigation.
